\section{Sliding Window Mode}
\label{sec:cses-sw-mode}

\problemstatement{Given an array of $n$ integers and a window size $k$,
output the mode of every contiguous window of size $k$ -- the value
occurring most often (if several tie, the smallest such value).}

\begin{patternbox}
The mode must survive both additions and \emph{removals}, so a single
``max frequency'' is not enough. Keep frequency buckets: for each frequency
$f$, an ordered set of the values that currently occur $f$ times. The mode
is the smallest value in the highest non-empty bucket.
\end{patternbox}

\subsection*{Frequency buckets of values}

Maintain $\textit{freq}[v]$ and $\textit{bucket}[f]$ = an ordered set of
values whose current frequency is exactly $f$, plus $\textit{maxFreq}$.
Adding a value $v$ moves it from $\textit{bucket}[\textit{freq}[v]]$ to
$\textit{bucket}[\textit{freq}[v]+1]$ and may raise $\textit{maxFreq}$.
Removing $v$ moves it down a bucket; if that empties the top bucket,
$\textit{maxFreq}$ decreases. The mode is the smallest value in
$\textit{bucket}[\textit{maxFreq}]$, read in $O(\log)$ from the ordered
set's begin.

\begin{ideabox}[Buckets Make the Mode Removal-Safe]
The trouble with modes under removal is that the top frequency can drop.
Grouping values by their exact frequency in ordered buckets lets a removal
adjust $\textit{maxFreq}$ in $O(1)$ and still expose the smallest tied
value instantly -- something a lone counter cannot do.
\end{ideabox}

\subsection*{Algorithm}

\begin{enumerate}
  \item Maintain $\textit{freq}$, $\textit{bucket}[f]$ (ordered set of
        values), and $\textit{maxFreq}$.
  \item On add/remove of $v$: erase $v$ from its old bucket, change its
        frequency, insert into the new bucket, and update
        $\textit{maxFreq}$.
  \item The mode is $\ast\textit{bucket}[\textit{maxFreq}].\text{begin}()$;
        output it per window.
\end{enumerate}

\subsection*{Dry run}

\dryrun{$a=[1,2,2,1,1]$, $k=3$.}

\begin{itemize}
  \item $[1,2,2]$: freqs $1{:}1,2{:}2$; maxFreq $2$, bucket$[2]=\{2\}$;
        mode $2$.
  \item Slide to $[2,2,1]$: still $2{:}2,1{:}1$; mode $2$.
  \item Slide to $[2,1,1]$: $1{:}2,2{:}1$; bucket$[2]=\{1\}$; mode $1$.
        Output $2,2,1$.
\end{itemize}

\subsection*{C++ implementation}

\begin{lstlisting}
#include <bits/stdc++.h>
using namespace std;

int main() {
    int n, k;
    cin >> n >> k;
    vector<int> a(n);
    for (int& x : a) cin >> x;

    unordered_map<int,int> freq;
    map<int, set<int>> bucket;                // frequency -> values
    int maxFreq = 0;

    auto change = [&](int v, int delta) {
        int f = freq[v];
        if (f > 0) bucket[f].erase(v);
        if (f > 0 && bucket[f].empty()) bucket.erase(f);
        f += delta;
        freq[v] = f;
        if (f > 0) bucket[f].insert(v);
        maxFreq = bucket.empty() ? 0 : bucket.rbegin()->first;
    };

    for (int i = 0; i < k; i++) change(a[i], +1);
    cout << *bucket[maxFreq].begin() << " ";
    for (int i = k; i < n; i++) {
        change(a[i], +1); change(a[i - k], -1);
        cout << *bucket[maxFreq].begin() << " ";
    }
    cout << "\n";
    return 0;
}
\end{lstlisting}

\begin{complexitybox}
\textbf{Time Complexity.} $O(n\log k)$ -- each add/remove touches ordered
sets. \textbf{Space Complexity.} $O(k)$.
\end{complexitybox}

\begin{takeawaybox}
Tracking the mode under removals needs frequency-to-values buckets so the
top frequency and its smallest value are both maintainable. The ``count of
counts'' idea -- indexing by frequency, not by value -- is the general trick
for order statistics on frequencies.
\end{takeawaybox}
